\documentclass[11pt,a4paper]{article}
% Preámbulo compartido por los dos manuales. Sobrio y sin colores: se imprime
% en blanco y negro en el mesón del taller.
\usepackage[T1]{fontenc}
\usepackage[utf8]{inputenc}
\usepackage[spanish,es-noquoting]{babel}
\usepackage[a4paper,margin=2.6cm]{geometry}
\usepackage{graphicx}
\usepackage{booktabs}
\usepackage{enumitem}
\usepackage{fancyhdr}
\usepackage[hidelinks]{hyperref}
\usepackage{titlesec}

\graphicspath{{imagenes/}}
\setlist{itemsep=2pt,topsep=4pt}
\titleformat{\section}{\large\bfseries}{\thesection}{0.6em}{}
\titleformat{\subsection}{\normalsize\bfseries}{\thesubsection}{0.6em}{}
\pagestyle{fancy}
\fancyhf{}
\fancyhead[L]{\small Lubri-Express — Sistema de Gestión}
\fancyhead[R]{\small\leftmark}
\fancyfoot[C]{\small\thepage}
\renewcommand{\headrulewidth}{0.4pt}

% Una captura de pantalla con su pie.
\newcommand{\pantalla}[2]{%
  \begin{figure}[htbp]
    \centering
    \includegraphics[width=\linewidth]{#1}
    \caption{#2}
  \end{figure}}

% Media captura, para los diálogos.
\newcommand{\dialogo}[2]{%
  \begin{figure}[htbp]
    \centering
    \includegraphics[width=0.62\linewidth]{#1}
    \caption{#2}
  \end{figure}}


\title{\vspace{-1.5cm}\textbf{Manual de usuario}\\[0.3cm]
       \large Sistema de Gestión de Taller, Inventario y Ventas}
\author{Lubri-Express}
\date{}

\begin{document}
\maketitle
\thispagestyle{empty}

\noindent Este manual describe el uso diario del sistema. La instalación y la
configuración están en el \emph{Manual técnico}.

\tableofcontents
\newpage

\section{Antes de empezar}

\subsection{Iniciar sesión}

El sistema pide usuario y contraseña cada vez que se abre. No hay sesión que
quede abierta de un día para otro: cada movimiento de inventario, cada venta y
cada orden quedan firmados con el nombre de quien los hizo, y esa firma es la
que permite después saber quién tocó qué.

\dialogo{login}{Inicio de sesión. Los errores se muestran bajo los campos.}

Si la contraseña no es correcta, o la cuenta está desactivada, el aviso aparece
en la misma ventana. Las contraseñas las cambia un administrador desde la
pestaña \textbf{Usuarios}.

\subsection{Qué puede hacer cada rol}

El sistema tiene tres roles. Los botones que un rol no puede usar aparecen
apagados, y al pasar el mouse por encima dicen a quién corresponden.

\begin{center}
\begin{tabular}{p{7.2cm}ccc}
\toprule
\textbf{Acción} & \textbf{Usuario} & \textbf{Supervisor} & \textbf{Admin.} \\
\midrule
Vender en mostrador & Sí & Sí & Sí \\
Abrir órdenes de trabajo & Sí & Sí & Sí \\
Clientes y vehículos & Sí & Sí & Sí \\
Consultar inventario y Kardex & Sí & Sí & Sí \\
Ver reportes & Sí & Sí & Sí \\
Registrar dinero y gastos en la caja del día & Sí & Sí & Sí \\
\midrule
Crear y editar productos & --- & Sí & Sí \\
Ingresar mercadería & --- & Sí & Sí \\
Ajustar stock por recuento & --- & Sí & Sí \\
Cargar inventario desde Excel & --- & Sí & Sí \\
Fijar el mínimo de una categoría & --- & Sí & Sí \\
Ver el precio de costo, en el listado y en el Kardex & --- & Sí & Sí \\
\midrule
Crear usuarios y asignar roles & --- & --- & Sí \\
\bottomrule
\end{tabular}
\end{center}

\subsection{Los precios son netos}

Todos los precios del catálogo —productos y servicios— se guardan \textbf{sin
IVA}. El impuesto se calcula al cobrar, se muestra desglosado y queda guardado
en la boleta o en la orden. Por eso, al crear un producto, el formulario pide el
precio neto y muestra al lado cuánto es con IVA.

\section{Inventario}

\pantalla{inventario}{Listado de productos, con el historial de movimientos del
producto seleccionado abajo.}

La tabla de arriba es el catálogo completo. Se busca escribiendo en la caja
superior: sirve el nombre, la marca o la categoría, sin importar tildes ni
mayúsculas. La casilla \textbf{Solo stock crítico} deja a la vista los productos
que están en su mínimo o por debajo.

Al pie de la pantalla, el resumen dice cuántos productos hay a la vista y
cuántos están bajo su mínimo. Para el supervisor y el administrador dice además
\textbf{cuánto vale a precio costo lo que está listado}: sin filtro es la bodega
entera y, con una búsqueda puesta, lo que vale esa marca o esa categoría. Va a
costo porque es la plata que está inmovilizada en la repisa, no lo que se va a
facturar por ella.

La \textbf{categoría} de un producto se elige de las que ya existen, y el
campo deja escribir una nueva cuando hace falta. Si se escribe una que ya
existe con otras mayúsculas o tildes, se guarda la que ya estaba: es lo que
evita terminar con \emph{Filtro aire}, \emph{Filtro Aire} y \emph{FILTRO
AIRE} como tres repisas distintas, que es como venía el sistema anterior.

Al elegir un producto, abajo aparece su \textbf{Kardex}: cada entrada, cada
salida por venta o por orden, con la fecha, la cantidad, el saldo que quedó y
quién lo hizo. Ese historial no se puede editar —es el registro auditable del
inventario— y por eso el stock nunca se escribe a mano.

\subsection{Recibir mercadería}

\dialogo{dialogo_ingreso}{Ingreso de mercadería: se arma la lista y se confirma
una sola vez.}

\begin{enumerate}
  \item \textbf{Ingresar Mercadería}.
  \item Buscar el producto. Al elegirlo, el campo \textbf{Costo unit.} se llena
        solo con lo que ese producto cuesta hoy.
  \item Poner la cantidad que llegó y, \textbf{si el proveedor cambió el
        precio}, corregir el costo. Después, \textbf{Añadir a la lista}. Se
        repite por cada producto de la factura.
  \item Revisar la columna \emph{Nuevo Stock}, que muestra cómo quedará cada uno.
  \item \textbf{Confirmar Ingreso}.
\end{enumerate}

Recién al confirmar se mueve el stock, y queda una entrada en el Kardex por cada
producto.

El costo que se escriba pasa a ser el del producto en el catálogo, y además
queda guardado en esa entrada del Kardex. Por eso, meses después, el historial
de un producto sigue diciendo a cuánto se compró cada vez. \textbf{El precio de
venta no cambia solo}: subir el precio es una decisión del taller y se hace en
la ficha del producto.

\subsection{Corregir el stock después de un recuento}

Cuando lo que hay en la repisa no coincide con lo que dice el sistema —una
merma, un error de digitación, un dato que venía mal del sistema antiguo— se
corrige contando.

\dialogo{dialogo_ajuste}{Ajuste por recuento: se escribe lo que hay de verdad.}

Se elige el producto, se aprieta \textbf{Ajustar stock} y se escribe el
\textbf{stock real contado}. El sistema calcula la diferencia y la registra como
un ajuste manual en el Kardex. El inventario queda corregido y, además, queda
dicho que se corrigió, cuándo y por quién.

\subsection{Cargar muchos productos desde Excel}

\dialogo{dialogo_excel}{Carga masiva: primero se ve qué entraría, y solo
entonces se carga.}

\begin{enumerate}
  \item \textbf{Cargar Excel} y después \textbf{Guardar plantilla\ldots}: el
        sistema entrega un archivo con las columnas que espera y una fila de
        ejemplo.
  \item Llenar la plantilla con LibreOffice o Excel.
  \item \textbf{Elegir planilla\ldots}: se ve la lista de lo que entraría, y
        cada fila dice si es un producto \emph{nuevo} o uno \emph{existente}.
  \item Si alguna fila tiene un problema, aparece abajo con su número de fila y
        el botón \textbf{Cargar} queda apagado. Se corrige la planilla y se
        vuelve a elegir.
\end{enumerate}

La carga es todo o nada, a propósito: si entrara una parte y después se repitiera
la planilla corregida, el stock de las filas que sí entraron quedaría sumado dos
veces. A un producto que ya existe solo se le suma el stock; sus precios no se
pisan desde una planilla.

\subsection{Stock mínimo por categoría}

\dialogo{dialogo_minimos}{El mínimo de toda una categoría, de una vez.}

El \emph{stock mínimo} es el umbral bajo el cual un producto aparece como
crítico. Fijarlo producto por producto es inviable con miles de ítems, así que
\textbf{Mínimo por categoría} lo aplica a todos los de una misma categoría. Un
producto puntual puede tener después su propio mínimo, editándolo.

\section{Ventas de mostrador}

\pantalla{ventas}{Punto de venta: catálogo a la izquierda, carrito a la derecha.}

\begin{enumerate}
  \item Buscar el producto y \textbf{Agregar al carrito} (o doble clic).
  \item Repetir por cada producto. La cantidad se ajusta al agregar.
  \item Si el cliente está registrado, elegirlo en \textbf{Cliente}; una venta
        sin cliente es normal y no exige nada.
  \item Escribir el \textbf{N.º de boleta}. Es obligatorio: es lo que amarra la
        venta del sistema con el documento que se le entrega al cliente.
  \item \textbf{Generar Venta}.
\end{enumerate}

El pie muestra el neto, el IVA y el total. Al cobrar, el stock baja solo y queda
la salida en el Kardex. Si entremedio alguien vendió lo último que quedaba, el
sistema lo avisa y no deja cobrar de más.

La pestaña \textbf{Historial de Ventas} lista lo cobrado; con doble clic se ve el
detalle de cada boleta.

\section{Clientes y vehículos}

\pantalla{clientes}{Clientes arriba; abajo, los vehículos del cliente elegido.}

Se busca por nombre, por RUT o por patente, sin importar tildes, puntos ni
guiones: \texttt{perez}, \texttt{Pérez} y \texttt{PEREZ} encuentran lo mismo, y
\texttt{12345678} encuentra al RUT \texttt{12.345.678-5}.

A las personas no se les pide RUT —en el mesón no lo dan y frena la atención—;
a las empresas sí, porque sin RUT no se les puede facturar.

\subsection{Vehículos}

\dialogo{dialogo_vehiculo}{Ficha del vehículo.}

Cada vehículo cuelga de un cliente y se identifica por su patente, que es única
en todo el sistema. Además de marca, modelo y año, la ficha guarda el tipo, la
versión, el VIN y el número de motor, que es lo que traía el sistema anterior.

\section{Órdenes de trabajo}

\pantalla{orden}{Orden en proceso. La tarjeta superior dice de qué vehículo se
trata y cuál fue su último servicio.}

\subsection{Abrir una orden}

\begin{enumerate}
  \item \textbf{Nueva Orden de Trabajo}.
  \item Elegir el cliente. Si no existe, \textbf{Crear Cliente} sin salir de ahí.
  \item Elegir el vehículo. Si es la primera vez que viene, \textbf{Añadir
        Vehículo}.
\end{enumerate}

Se abre la mesa de trabajo, con una tarjeta arriba que muestra la patente, el
vehículo, el cliente y \textbf{el último servicio registrado con su
kilometraje}: es contra ese número que se compara el de hoy.

\subsection{Cargar lo que se usó}

Los \textbf{productos} (aceites, filtros, insumos) descuentan stock; los
\textbf{servicios} (mano de obra) no, y por eso van en su propio buscador. Se
agregan uno por uno, con su cantidad.

Si se sabe el nombre, se escribe en el campo \textbf{Producto} y la lista se va
reduciendo. Cuando hay varios que se llaman casi igual ---dos filtros de aire
del mismo modelo, uno de cada marca--- conviene el botón \textbf{Buscar\ldots}:
abre una ventana con todo el catálogo y, al lado de cada nombre, la marca, la
categoría, \textbf{cuántos quedan}, \textbf{en qué repisa están} y el precio. El
stock bajo el mínimo sale en rojo. Se elige con doble clic y el producto queda
puesto en el campo, listo para agregar con su cantidad.

\dialogo{dialogo_buscar_producto}{El buscador de productos: stock y ubicación al
lado de cada nombre.}

El \textbf{kilometraje de ingreso} es obligatorio: es con lo que se calcula el
próximo servicio. Si se anota un kilometraje menor que el del servicio anterior,
el sistema pregunta antes de guardar —suele ser un dedo cambiado, pero a veces
es un tablero reemplazado.

\subsection{Dejar la orden abierta (el auto se queda)}

\pantalla{ordenes_abiertas}{Los autos que siguen en el taller, cada uno con su
orden esperando.}

No todos los autos se entregan el mismo día. Si el trabajo sigue mañana, o falta
un repuesto, el botón \textbf{Dejar abierta} guarda la orden con el auto todavía
adentro. Se pueden tener \textbf{varias órdenes abiertas a la vez}, una por auto.

Aparecen en la pestaña \textbf{Órdenes Abiertas}, que muestra entre paréntesis
cuántas hay. Ahí se elige una y:

\begin{itemize}
  \item \textbf{Retomar} la trae de vuelta a la mesa de trabajo tal como quedó,
        con sus insumos cargados. Se le agrega lo que falte y se entrega con
        \textbf{Guardar y entregar}.
  \item \textbf{Anular} cierra la orden sin entregarla.
\end{itemize}

\textbf{Lo importante del stock.} Una orden abierta \emph{ya descontó} sus
repuestos: es correcto, porque el mecánico los sacó de la repisa cuando los usó.
Por eso, si se le quita una línea al retomarla, o si se anula la orden entera,
\textbf{ese stock vuelve solo a la bodega} y queda anotado en el Kardex diciendo
de qué orden volvió. No hay que corregir nada a mano.

Una orden anulada no desaparece: queda en el historial marcada como tal, porque
el auto sí entró al taller. Los reportes no la cuentan.

\subsection{Descuento, folio y pago}

El descuento se aplica \textbf{en porcentaje o en monto fijo}, nunca los dos a
la vez: se elige cuál en la lista y se escribe el valor al lado. El IVA se
calcula sobre lo que efectivamente se cobra, ya descontado.

El \textbf{folio de Mercado Público} se llena solo en las órdenes
institucionales; es lo que después permite separarlas de los clientes normales
en el historial.

El pago se anota al cerrar la orden. Si el cliente paga todo, se marca la
casilla \textbf{Pagada}. Si paga una parte —un abono al dejar el auto y el resto
al retirarlo— se escribe lo que dejó en \textbf{Abono}, que no acepta más que el
total de la orden. El saldo se cobra después, desde el detalle de la orden en el
historial; no hay que anotarlo en ninguna parte ni llevar la cuenta a mano.

\subsection{Historial y PDF para el cliente}

\pantalla{historial_ordenes}{Historial de órdenes, con el filtro de Mercado
Público.}

El historial se busca por número de orden, patente o cliente, y el filtro de la
derecha separa \emph{Mercado Público} de \emph{Clientes}. Con doble clic se abre
el detalle de la orden.

La columna \textbf{Estado} dice cómo va el cobro: \emph{Pagada},
\emph{No pagada}, o \emph{Abonada} cuando quedó a medio pagar —poniendo el
cursor encima aparece cuánto abonó y cuánto debe.

\dialogo{dialogo_detalle_orden}{Detalle de una orden guardada.}

Desde ahí, \textbf{Exportar PDF} genera el documento para el cliente, con el
desglose de insumos y servicios, el descuento, el IVA y el total. Se guarda en
\texttt{Documentos\textbackslash Lubri-Express\textbackslash Órdenes}.

\textbf{Registrar pago} anota lo que el cliente paga ahora. Viene propuesto el
saldo completo, que es el caso de todos los días —viene a buscar el auto y paga
lo que debe—, y se puede escribir menos si deja otro abono. Los pagos se suman y
no se corrigen: la orden pasa a \emph{Pagada} sola, cuando lo abonado alcanza el
total. Mientras quede saldo, el PDF del cliente muestra lo abonado y lo que
falta.

\subsection{Avisar al cliente por WhatsApp}

El botón \textbf{Avisar por WhatsApp} está en dos lugares, y los dos son la
orden: en la tarjeta del vehículo mientras la orden está abierta, para contar
cómo va el trabajo, y en el detalle de una orden guardada, que es desde donde se
avisa que el auto está listo —ahí el mensaje sale citando el número de la OT.

Se elige una de las cuatro plantillas —\emph{Listo para retiro}, \emph{Trabajo
en curso}, \emph{Requiere su autorización}, \emph{Esperando repuesto}— y el
mensaje aparece redactado, con el nombre del cliente y su vehículo.
\textbf{Se puede editar antes de mandarlo}: la plantilla está para no escribir lo
mismo veinte veces al día, no para mandarla a ciegas.

\textbf{Abrir en WhatsApp} abre WhatsApp —el de escritorio si está instalado, o
WhatsApp Web— con el mensaje ya escrito en la conversación. \emph{No lo envía}:
el último clic lo da quien atiende, que es quien sabe si el auto está listo de
verdad.

El botón está apagado si el cliente no tiene registrado un celular chileno; un
teléfono fijo no sirve, WhatsApp abriría una conversación que no existe. El
número se corrige en la ficha del cliente, en \textbf{Clientes}.

\section{Caja}

\pantalla{caja}{La caja del día: lo agregado, lo cobrado, los gastos y con
cuánto debería quedar el cajón.}

Es la \textbf{caja chica del día}: sirve para saber con cuánta plata debería
quedar el cajón de la oficina al cerrar.

\textbf{No hay que abrirla ni cerrarla.} La caja de un día es simplemente la de
esa fecha: al día siguiente empieza una nueva, sola. Arriba a la izquierda se
elige el día; se puede mirar cualquier día anterior, pero no uno futuro.

Las cuatro cifras son:

\begin{itemize}
  \item \textbf{Agregado}: el efectivo que se puso en el cajón ese día, anotado
        a mano ---el de la mañana para dar vuelto, por ejemplo---.
  \item \textbf{Ingresos}: lo cobrado ese día por ventas de mostrador y por
        abonos de órdenes. \textbf{Esto entra solo}, no hay que anotarlo.
  \item \textbf{Egresos}: los gastos anotados a mano.
  \item \textbf{Balance}: agregado más ingresos, menos egresos. Es lo que
        debería haber en el cajón.
\end{itemize}

\subsection{Agregar dinero y anotar un gasto}

\textbf{Agregar dinero} es para el efectivo que se pone en el cajón;
\textbf{Agregar gasto} para lo que sale de ahí (bencina, un flete, el almuerzo).
Los dos piden el monto y \textbf{un motivo, que es obligatorio}: un gasto sin
motivo no sirve para nada el día que hay que cuadrar.

Al agregar dinero, el sistema propone el saldo con que cerró el último día que
tuvo movimiento ---un lunes propone el del viernes---. Casi siempre es la plata
que quedó en el cajón, y así no hay que acordarse de la cifra.

\subsection{Si se anotó algo mal}

Las líneas de la caja \textbf{no se borran ni se editan}. Se elige la línea
equivocada y se aprieta \textbf{Anular}: el sistema anota el movimiento
contrario. Las dos líneas quedan a la vista, en gris, y dejan de contar para las
cifras del día.

Es la misma regla del Kardex: lo que pasó por caja queda registrado, y quién lo
anotó también.

\subsection{Guardar la caja del día}

\textbf{Guardar caja en PDF} deja un archivo en
\texttt{Documentos\textbackslash Lubri-Express\textbackslash Caja\textbackslash}
con las cuatro cifras y el detalle de lo anotado a mano. Sirve para imprimirlo o
archivarlo al cierre.

\section{Reportes}

\pantalla{reportes}{Ingresos del período, con su gráfico.}

Arriba se elige el rango de fechas. Hay cuatro reportes:

\begin{itemize}
  \item \textbf{Ingresos por período}: ventas y órdenes por día —o por mes, si
        el rango es largo—, con neto, IVA y total.
  \item \textbf{Ventas por producto}: cuánto se vendió de cada cosa.
  \item \textbf{Por usuario}: cuánto movió cada persona.
  \item \textbf{Reabastecimiento}: qué hay que reponer.
\end{itemize}

\pantalla{reabastecimiento}{Reabastecimiento: los productos en su mínimo o por
debajo, ordenados por lo que falta.}

\textbf{Exportar a Excel} guarda el reporte que está a la vista en
\texttt{Documentos\textbackslash Lubri-Express\textbackslash Reportes}, y
\textbf{Abrir carpeta de reportes} abre esa carpeta.

\medskip
\noindent\emph{Nota:} las órdenes traídas del sistema anterior se cargaron con
su total, pero sin el detalle de qué insumos llevaron —el sistema viejo no lo
guardaba—. Por eso \emph{Ventas por producto} refleja lo vendido desde que se
usa este sistema, mientras que \emph{Ingresos por período} sí cubre todo el
historial.

\section{Usuarios}

\pantalla{usuarios}{Administración de usuarios. Solo la ve el administrador.}

\textbf{Nuevo usuario} pide nombre, usuario, rol y contraseña. Al editar, dejar
la contraseña en blanco la conserva; escribir una nueva la reemplaza.

Un usuario \textbf{desactivado} no puede entrar, pero su nombre sigue firmando
todo lo que hizo: por eso se desactiva en vez de borrarse. Un administrador no
puede desactivarse a sí mismo —el sistema quedaría sin nadie que pueda
deshacerlo—.

\section{Preguntas frecuentes}

\subsection*{El producto no aparece al vender}
Puede estar desactivado. En Inventario se ve en gris y su fila lo dice.

\subsection*{Aparecen muchos productos como ``stock crítico''}
El mínimo empieza en cero para todo lo que vino del sistema anterior, así que
cualquier producto sin stock figura como crítico. Se corrige fijando los mínimos
por categoría (sección 2.4).

\subsection*{Guardé una orden y el auto todavía no estaba listo}

Para eso es \textbf{Dejar abierta}. Si ya la guardaste con \textbf{Guardar y
entregar}, no se puede reabrir: se hace una orden nueva por lo que falte.

\subsection*{Me equivoqué en una orden ya guardada}
Las órdenes no se editan: son el registro de lo que se hizo. Si hay que corregir
el inventario, se hace con un ajuste por recuento, que deja constancia. Lo único
que se le agrega a una orden guardada es un pago.

\subsection*{El balance de la caja no coincide con el cajón}

Primero, revisar que estén anotados \emph{todos} los gastos del día: es lo que
más se olvida. Después, tener presente que \textbf{la caja cuenta todo lo
cobrado}, sin distinguir efectivo de tarjeta o transferencia; si ese día hubo
pagos que no entraron en efectivo, el cajón va a tener menos que el balance.

Y si un día antiguo aparece con \emph{Ingresos} en cero: las órdenes que vienen
del sistema anterior no traen el detalle de cómo se pagaron, así que no suman.

\subsection*{El sistema dice que no se pudo conectar a la base de datos}
Es un problema de instalación, no de uso. Está en el Manual técnico.

\end{document}
